% !PRINTED 545
\weberpage{545}
But we have
\[
\begin{aligned}
K(-3) &= \tfrac{1}{3}, & K(-84) &= 4, & K(-20) &= 2, & K(-35) &= 2, \\
K(+140) &= 4, & K(5) &= 1, & K(21) &= 2, & K(12) &= 2^{1)},
\end{aligned}
\]
\begingroup\renewcommand{\thefootnote}{1)}\footnotetext[1]{According to the notation of \gesperrt{Gauss} the class numbers of the second kind are to be taken.}\endgroup
as one finds from the tables of \gesperrt{Gauss}, or here also easily by direct enumeration.

Further, according to the tables of \gesperrt{Legendre} or of \gesperrt{Degen}:
\[
\begin{aligned}
\log\left(\frac{T+\sqrt{140}\,U}{2}\right) &= \log(6+\sqrt{35}) = \log\frac{(\sqrt{5}+\sqrt{7})^2}{2}, \\
\log\left(\frac{T+\sqrt{5}\,U}{2}\right) &= \log\left(\frac{1+\sqrt{5}}{2}\right)^2 = \tfrac{2}{3}\log(2+\sqrt{5}), \\
\log\frac{T+\sqrt{21}\,U}{2} &= \log\frac{5+\sqrt{21}}{2} = \tfrac{1}{3}\log(55+12\sqrt{21}) \\
&= \log\left(\frac{\sqrt{3}+\sqrt{7}}{2}\right)^2, \\
\log\left(\frac{T+\sqrt{12}\,U}{2}\right) &= \log(2+\sqrt{3}),
\end{aligned}
\]
and from this one obtains:
\[
\left(\frac{f(\sqrt{-105})}{\sqrt[4]{2}}\right)^{12} = (2+\sqrt{5})^2(55+12\sqrt{21})(6+\sqrt{35})(2+\sqrt{3})^3,
\]
or, what means the same thing:
\[
\sqrt{2}^{13}f(\sqrt{-105})^6 = (1+\sqrt{5})^3(\sqrt{3}+\sqrt{7})^3(\sqrt{5}+\sqrt{7})(1+\sqrt{3})^3.
\]

\begin{center}
\textbf{\S\,144. Calculation of some further class invariants.}
\end{center}

Next after the discriminants for which only \gesperrt{one} class of forms occurs in each genus, which we considered in the preceding section, the simplest results are given by those which contain two classes of forms in each genus, and among these again by those for which only two genera occur. The class invariants for these discriminants $-4m$ are the roots of quadratic equations whose coefficients contain only one square root. We again assume fundamental discriminants, and obtain from the table of \gesperrt{Gauss} the following values of $m$:
%CONT>>
